% Appendix: formal definition of calibration, Murphy/Bröcker decomposition, and post-hoc calibration (Exp 41).
% One sentence per line (Trigger Rule 13); strictly anonymized.

\section{Calibration: Formal Definition}
\label{app:calibration}

This appendix defines calibration, records the associated proper-score decomposition, and states why correlation does not see it.
The main text writes $x$ for the input, $r$ for the rating, $p(r\mid x)$ for the human rating distribution, and $q(r\mid x)$ for the judge's forecast.
Here we use the equivalent notation $c$ for the context (the same object as $x$), $Y$ for the human rating, $p^*(\cdot\mid c)$ for the human rating distribution of a stated rater pool and rubric, and $\hat q(\cdot\mid c)$ for the judge's forecast, with $C$ the random context.

\subsection{Definition}
Let $Y\in\{1,\dots,K\}$ and let $\hat q(\cdot\mid c)$ be a forecast distribution on these $K$ labels.
The forecast is \emph{calibrated} if, whenever it issues a particular distribution $q$, the human rating has that distribution:
\begin{equation}
\label{eq:calibration}
\Pr\bigl(Y=k \mid \hat q(\cdot\mid C)=q\bigr)
=
q_k
\qquad\text{for all }q\text{ and all }k.
\end{equation}
In the binary case $Y\in\{0,1\}$, writing $\hat p(c)=\hat q(Y=1\mid c)$, this reduces to
\begin{equation}
\label{eq:calibration-binary}
\Pr\bigl(Y=1 \mid \hat p(C)=p\bigr)
=
p
\qquad\text{for all }p\in[0,1].
\end{equation}
Intuitively, a calibrated judge's stated probabilities are the frequencies with which the corresponding human ratings actually occur: when the judge says $q$, the humans do $q$.

Calibration is a statement about $\Pr(Y\mid\hat q)$, not about whether $\hat q(\cdot\mid c)$ equals $p^*(\cdot\mid c)$ for each $c$.
A forecast can be calibrated yet coarse (many contexts share the same $q$), and an uncalibrated forecast can still discriminate contexts well.

\subsection{Proper-Score Decomposition}
Write $S(q,y)$ for a scoring rule in the loss orientation of the main text (lower is better).
If $S$ is strictly proper, it induces a divergence $d(q,p)=\mathbb{E}_{Y\sim p}[S(q,Y)]-\mathbb{E}_{Y\sim p}[S(p,Y)]\ge 0$, with equality if and only if $q=p$~\citep{gneiting2007strictly}.
Let $p_{\hat q}(\cdot)=\Pr(Y\in\cdot\mid\hat q)$ be the recalibrated conditional law and let $\bar p$ be the marginal (base-rate) law of $Y$.
The expected score then decomposes as
\begin{equation}
\label{eq:proper-decomposition}
\mathbb{E}\bigl[S(\hat q,Y)\bigr]
=
\mathbb{E}\bigl[S(\bar p,Y)\bigr]
-
\mathbb{E}\bigl[d(\bar p,p_{\hat q})\bigr]
+
\mathbb{E}\bigl[d(\hat q,p_{\hat q})\bigr].
\end{equation}
The three terms are, in order, \emph{uncertainty} (the score of the base-rate forecast), \emph{resolution} (how much the recalibrated law varies across issued forecasts), and \emph{calibration error}
\begin{equation}
\label{eq:calerror}
\mathrm{CAL}
=
\mathbb{E}\bigl[d\bigl(\hat q,\Pr(Y\mid\hat q)\bigr)\bigr].
\end{equation}
The calibration error is zero if and only if $\hat q=\Pr(Y\mid\hat q)$ almost surely, which is exactly calibration~\eqref{eq:calibration}.
For the Brier score this is the vector partition of \citet{murphy1973vector}.
\citet{brocker2009reliability} extends the decomposition to strictly proper scores on finite-valued targets;~\citet{gneiting2007strictly} supply the associated divergence used in~\eqref{eq:proper-decomposition}.
Forecast verification calls the third term \emph{reliability}; we call it calibration error throughout so that it is not confused with psychometric reliability (\appref{app:reliability}{the reliability appendix}).

The identity~\eqref{eq:proper-decomposition} holds for any score $S$, because it uses only the linearity of expectation and the tower property; propriety is what makes the second and third terms nonnegative, and strict propriety is what makes the third zero only under calibration.
The conditioning is on the issued forecast value.
Conditioning on a finer partition, for example the item, leaves the sum unchanged and moves part of the calibration error into resolution, so the finer key never lowers the third term.
Status of the decomposition: for squared error with a scalar score and for the multiclass Brier score it is machine-checked in Lean~4 for finite laws (\texttt{Point2.FinLaw.murphy}, \texttt{murphy\_\allowbreak scalar}, \texttt{brier\_\allowbreak murphy}, \texttt{reliability\_\allowbreak zero\_\allowbreak iff}, \texttt{brier\_\allowbreak reliability\_\allowbreak zero\_\allowbreak iff} in \texttt{docs/\allowbreak notes/\allowbreak point2\_\allowbreak support\_\allowbreak 27Nov2026/\allowbreak lean/\allowbreak Point2/\allowbreak Murphy.lean}; the Lean names say \texttt{reliability} for the calibration error; Lean 4.36.0, Mathlib revision \texttt{fc4f59e}, zero \texttt{sorry}, standard axioms only).
For cross-entropy and the ranked probability score it is numerically checked and the Lean proof is pending: a seeded script checks it on $300$ random finite laws per score (exact fractions for the Brier score and RPS, tolerance $10^{-12}$ for cross-entropy), together with the claims that must fail (\texttt{proper\_\allowbreak score\_\allowbreak decomposition\_\allowbreak numerics\_\allowbreak 28Nov2026.py}, $25$ of $25$ checks pass).

\subsection{Why Correlation Misses Calibration}
As in the main text and \appref{app:concordance}{the concordance appendix}, squared Pearson correlation admits the identity
\begin{equation}
\label{eq:r-squared-cal}
r^2
=
1
-
\frac{\mathrm{MSE}(\text{best affine fit})}{\mathrm{Var}(Y)}.
\end{equation}
Correlation therefore equals one minus standardized squared error after an optimal affine refit of the judge's scalar to $Y$.
Any miscalibration that an affine map can remove is invisible, and spread about the conditional mean is never scored.
In the language of~\eqref{eq:proper-decomposition}, correlation keeps only the part of the resolution that the affine refit recovers.
For squared error and a scalar score $S$, with $g(s)=\mathbb{E}[Y\mid S=s]$ and $S'$ the best affine refit of $S$ to $Y$, the identity is $r^2\,\mathrm{Var}(Y)=\mathrm{Var}\,g(S)-\mathbb{E}\bigl(g(S)-S'\bigr)^2$ (Lean: \texttt{Point2.FinLaw.r\_\allowbreak sq\_\allowbreak resolution}).
So $r^2\,\mathrm{Var}(Y)$ never exceeds the resolution and equals it exactly when the calibration curve $g$ is affine in $S$.
An affine relabelling of the score leaves $r$ and $g$ unchanged and changes the calibration error (Lean: \texttt{condE\_\allowbreak affine}).
Correlation is therefore blind to miscalibration that an affine map can remove, not to every miscalibration: an informative but nonmonotone score can have $r=0$ with maximal resolution ($Y=S^2$), and a strictly increasing nonlinear relabelling can change Pearson's $r$.
These last two statements are numerically checked on exact finite examples (the script of the previous subsection, checks 8f and 9g).
Lin's concordance identity~\eqref{eq:lin-concordance} makes the same split visible as bias and scale terms that $r$ does not constrain.

\subsection{Post-Hoc Calibration}
Given a frozen forecast $\hat q$ and a held-out validation sample $\{(c_i,y_i)\}_{i=1}^n$, a post-hoc map $g$ is chosen by empirical proper-score risk
\begin{equation}
\label{eq:posthoc-cal}
\hat g
=
\arg\min_{g\in\mathcal{G}}
\frac{1}{n}\sum_{i=1}^{n} S\bigl(g(\hat q(\cdot\mid c_i)),y_i\bigr),
\end{equation}
where $\mathcal{G}$ is a restricted family such as temperature scaling, Platt scaling, or isotonic regression.
The aim is to reduce~\eqref{eq:calerror} without retraining the judge.
A map removes at most the calibration error: for every map $g$, $\mathbb{E}S(\hat q,Y)-\mathbb{E}S(g(\hat q),Y)\le\mathrm{CAL}$, with equality at $g(q)=\Pr(Y\mid\hat q=q)$ (numerically checked, check 6 of the script above).
The calibration errors reported for real judges below are therefore lower bounds, within the map family, the folds and the sample used.
A strictly increasing $g$ of a scalar score leaves the ranking of items unchanged; isotonic regression is only nondecreasing and can introduce ties.

This is the map used in Experiment~41.
On seven real multi-annotator rating sets re-binned to two, three and five classes, a cross-entropy-trained linear head beat a Pearson-trained one by up to $0.1282$ normalized ranked probability score (RPS) on held-out items (fixed ridge penalty $0.01$; the largest gap is on a view where neither head beat the constant forecast).
After each head's scalar passed through its own validation-fitted calibration map of this form, the largest five-seed mean difference was $0.0024$ (single seeds up to $0.0031$).
As skill over the base rate (Experiment~53), the Pearson-trained head's raw skill is negative on $15$ of $16$ views and the calibrated skill difference is $-0.0003$ to $+0.0098$.
The gap that correlation could not see was almost entirely calibration error in the sense of~\eqref{eq:calerror}.
Frozen frontier judges show a similar calibration gap (Experiment~46, interim): on DICES-350, four of five score worse than the base rate in normalized RPS before calibration (skill $-0.177$ to $-0.096$, each 95\% item-cluster bootstrap interval below $0$) at $80$ to $82$ percent top-1 agreement with the human majority, and every judge's skill is positive on both DICES sets after a per-judge calibration map fitted on held-out folds.
Experiment~57 splits the normalized-RPS risk of the same five judges into the terms of~\eqref{eq:proper-decomposition} (calibration map fitted on other folds, five group folds, each term divided by the base-rate risk so that it is in skill units, higher resolution is better, lower calibration error is better).
The calibration error ranges from $0.037$ to $0.231$ on DICES-350 and from $0.029$ to $0.207$ on DICES-990, whereas the resolution ranges only from $0.050$ to $0.072$ and from $0.090$ to $0.120$; the between-judge standard deviation is $0.077$ against $0.010$ on DICES-350 and $0.072$ against $0.012$ on DICES-990.
One judge shows the split: on DICES-350, GPT-6 Astra has Pearson correlation $0.588$ with the human item mean (second of five), raw skill $-0.143$ (fourth of five; $95\%$ interval $-0.188$ to $-0.100$), calibration error $0.215$ and the highest resolution ($0.072$, a near tie with $0.069$ for Claude Opus 5.5).
We had pre-registered that the correlation ranking of the five judges tracks their resolution ranking better than their raw-skill ranking does in at least four of five rating sets; that prediction was falsified (it held in one set, DICES-350), so on these judges the disagreement between correlation and a proper score concerns level more than order (Experiment~57).
